% Generated draft: requires mathematical and source review.
\documentclass[11pt]{article}
\usepackage[margin=25mm]{geometry}
\usepackage{fontspec}
\setmainfont{Latin Modern Roman}
\usepackage{amsmath,amssymb,mathtools,unicode-math}
\setmathfont{Latin Modern Math}
\usepackage{enumitem,needspace,xurl,hyperref}
\setlength{\parindent}{0pt}
\setlength{\parskip}{5pt}
\setlength{\emergencystretch}{3em}
\newcommand{\N}{\mathbb{N}}
\newcommand{\Z}{\mathbb{Z}}
\newcommand{\Q}{\mathbb{Q}}
\newcommand{\R}{\mathbb{R}}
\newcommand{\C}{\mathbb{C}}
\newcommand{\F}{\mathbb{F}}
\newcommand{\A}{\mathbb{A}}
\newcommand{\PP}{\mathbb{P}}
\newcommand{\E}{\mathbb{E}}
\newcommand{\eps}{\varepsilon}
\newcommand{\dd}{\,\mathrm{d}}
\begin{document}
\section*{UnsolvedMath ID: 800101}
\textbf{Source problem code:} catalogue entry
\subsection{Definitions and mathematical statement}
Let \(G\) be a group and let \(\operatorname{Homeo}^{+}(\R^{2})\) denote the group of orientation-preserving homeomorphisms of the plane. Consider a homomorphism
\[
\Phi\colon G\longrightarrow\operatorname{Homeo}^{+}(\R^{2})
\]
preserving two foliations \(\mathcal F^{+}\) and \(\mathcal F^{-}\). The source allows possible prong singularities and describes the foliations as transverse; their precise admissible singularities and the meaning of transversality at singular points are not specified.

The source states that this action extends to an associated circle at infinity \(S^{1}_{\infty}(\mathcal F^{+},\mathcal F^{-})\), producing an induced action
\[
\varphi\colon G\longrightarrow
\operatorname{Homeo}^{+}\!\left(S^{1}_{\infty}(\mathcal F^{+},\mathcal F^{-})\right).
\]
This extension assertion is supplied background attributed to the reference key \texttt{bonatti2024action}; its precise hypotheses and the circle construction are not included in the record.

\textbf{Question.} Can smoothability of the plane action \(\Phi\) be detected from its induced circle action \(\varphi\)? In particular, what condition on the induced action characterizes smoothability of the original plane action?

For a specified differentiability class \(C^{r}\), a possible precise interpretation of smoothability is the existence of a homeomorphism \(h\colon\R^{2}\to\R^{2}\) such that
\[
 h\circ\Phi(g)\circ h^{-1}\in\operatorname{Diff}^{r}_{+}(\R^{2})
 \qquad\text{for every }g\in G,
\]
where \(\operatorname{Diff}^{r}_{+}(\R^{2})\) denotes the orientation-preserving \(C^{r}\) diffeomorphisms of the plane. Here the same \(h\) must work for every group element. This interpretation is conditional: the source does not specify the differentiability class, the intended notion of smoothability, whether the foliations must also become smooth, or what equivalence and additional data are allowed when inspecting \(\varphi\). These choices must be resolved before the question has a unique fully precise formulation.
\subsection{Short English statement}
An orientation-preserving action on the plane preserves two transverse foliations, possibly with prong singularities, and the source assigns it an induced action on a circle at infinity. The question asks whether that circle action reveals whether the original plane action can be made smooth. The meaning of smoothability and the exact hypotheses for the boundary construction remain unspecified.
\subsection{Sources}
\begin{itemize}
\item Local UnsolvedMath registry, permanent ID \(800101\), exact source problem code \texttt{catalogue entry}, title \emph{Foliations and Diffeomorphism Groups --- Question 42}. The supplied registry labels the problem \emph{open}; this label is not an independently verified mathematical or literature assessment.
\item Nariman and Yazdi, \emph{Problem List on Foliations and Diffeomorphism Groups} (2025), Question 42, \url{https://arxiv.org/abs/2503.18274}. The record attributes the problem to Nariman and Yazdi, with proposed year 2025. Supplied provenance records access on 2026-10-16 and extraction method \texttt{arxiv\_latex}; these are supplied provenance claims, not a claim of independent source access.
\item The supplied statement attributes the extension to the circle at infinity to \texttt{bonatti2024action}. No full bibliographic entry, title, author list, or URL for that reference is supplied, so its content and hypotheses cannot be checked from this record.
\end{itemize}
\end{document}
